\SourceLeaf{405}%
where \(c\) is the difference between the radii of the first and second circles and \(c'\) the difference between the radii of the first and third.

We have the following equations
\[
\begin{aligned}
xx+yy&=zz\\
(x-a)^2+(y-b)^2&=(z-c)^2\\
(x-a')^2+(y-b')^2&=(z-c')^2
\end{aligned}
\]
If we set
\[
\begin{aligned}x&=z\cos\varphi\\y&=z\sin\varphi\end{aligned}
\]
we obtain from the preceding equations
\[
\begin{aligned}
aa+bb-cc&=2az\cos\varphi+2bz\sin\varphi-2cz\\
a'a'+b'b'-c'c'&=2a'z\cos\varphi+2b'z\sin\varphi-2c'z
\end{aligned}
\]
Divide the first equation by \(a\cos\varphi+b\sin\varphi-c\) and the second by \(a'\cos\varphi+b'\sin\varphi-c'\). Subtracting one from the other then gives
\[
\frac{a'a'+b'b'-c'c'}{a'\cos\varphi+b'\sin\varphi-c'}-\frac{aa+bb-cc}{a\cos\varphi+b\sin\varphi-c}=0
\]
or, if we write
\[
\begin{aligned}a'a'+b'b'-c'c'&=A'\\aa+bb-cc&=A\end{aligned}
\]
this becomes
\[
(A'a-Aa')\cos\varphi+(A'b-Ab')\sin\varphi=A'c-Ac'
\]
If we now set
\[
\begin{aligned}A'a-Aa'&=R\cos M\\A'b-Ab'&=R\sin M\\A'c-Ac'&=N\end{aligned}
\]
our equation becomes
\[
R\cos(\varphi-M)=N
\]
where everything except \(\varphi\) is given. From this we get two values for \(\varphi\), and we must choose between them according to how the required circle is to touch the given circles.
